\begin{table}[htbp]
\centering
\caption{Revenue-per-Capita Border Gap by State Market Age}
\label{tab:revenue_percapita_eventstudy}
\begin{tabular}{lcccc}
\toprule
Years since state's own retail launch & Border $\times$ Year coef. & SE (county-cl.) & $p$-value & Border obs.\\
\midrule
0 & $-0.421$ & $0.360$ & $0.243$ & 7 \\
1 & $+0.142$ & $0.226$ & $0.531$ & 7 \\
2 & $-0.287$ & $0.357$ & $0.422$ & 9 \\
3 & $-0.353$ & $0.343$ & $0.304$ & 11 \\
4 & $-0.168$ & $0.282$ & $0.551$ & 11 \\
5 & $-0.049$ & $0.279$ & $0.862$ & 11 \\
6 & $+0.085$ & $0.265$ & $0.750$ & 12 \\
7 & $+0.153$ & $0.288$ & $0.596$ & 12 \\
8 & $+0.304$ & $0.267$ & $0.255$ & 12 \\
9+ & $+0.334$ & $0.291$ & $0.251$ & 30 \\
\bottomrule
\end{tabular}
\begin{minipage}{0.95\textwidth}
\vspace{0.3em}
\footnotesize
\textit{Notes:} Outcome is log(revenue per capita); same specification and panel as
Table~\ref{tab:revenue_perstore_eventstudy}, $n=815$. No coefficient approaches significance at any
market age -- the per-capita null (Table~\ref{tab:revenue_border70}) holds throughout a state's market
lifecycle, not merely on average. Confidence intervals are wide and overlap zero in every bin, in sharp
contrast to Figure~\ref{fig:revenue_perstore_eventstudy}'s per-store premium, which is precisely the
contrast the license-cap mechanism in Section~\ref{sec:revenue_perstore} predicts: if the number of stores
cannot expand to meet elevated border demand, that demand should not show up when revenue is divided by
population, only when it is divided by the (capped) number of sellers.
\end{minipage}
\end{table}
